\chapter{Addresses and memory}
\MARSChapterHeading{Addresses and memory}
\label{revision:addresses-and-memory}

The byte inspector kept its working values in registers. A larger program also
needs storage that can hold many values and remain available between operations.
Memory provides that storage. To use it correctly, we must distinguish a location
from the contents of that location, and specify how many bytes an instruction
reads or writes.

Our program will write one RGB color into memory, replace its blue component,
and read the result back. MARS's Bitmap Display can show the color, but the memory
bytes will be our precise evidence. We will also keep a recognizable value beside
the pixel and check that changing one byte does not damage its neighbor.

\section{A location has an address}

MARS presents memory as individually addressed bytes. An \emph{address} identifies
one of those byte locations. Consecutive addresses differ by one, even when our
program groups several bytes into a larger value. A 32-bit word occupies four
consecutive bytes. If its first byte is at address \texttt{0x10010000}, its other
bytes are at \texttt{0x10010001}, \texttt{0x10010002} and \texttt{0x10010003}.
The next nonoverlapping word can begin at \texttt{0x10010004}.

A label gives us a name for a location without repeating its numerical address
throughout the source. In the declaration \texttt{pixel: .word 0}, the label
\texttt{pixel} names the first byte of a four-byte object, and \texttt{.word 0}
reserves those bytes with an initial value of zero. On this MIPS64 target,
\texttt{.word} still describes four bytes. The general-purpose registers have
eight bytes each; that does not change the directive's storage size.

The instruction \texttt{la \$t0, pixel} puts the label's address in
\texttt{\$t0}. It does not read the word stored at that address. We can think of
the register as holding a pointer: a value our program uses to locate an object.
An address in a register is still a bit pattern. The subsequent instruction
determines whether to use it as a location or as an arithmetic operand.

For the complete program below, select MARS's Default memory layout. Its first
static data word is at \texttt{0x10010000}. Keeping \texttt{pixel} first lets the
display header and the label refer to the same location. In ordinary code, using
the label avoids depending on that numerical address. Here, the display interface
also needs a known base, so the data order is part of the example's contract.

\section{Moving a word between registers and memory}

We will put a color value in \texttt{\$t1} and its destination address in
\texttt{\$t0}. The instruction \texttt{sw \$t1, 0(\$t0)} stores the low 32 bits
of \texttt{\$t1} into four bytes beginning at that address. The letters
\texttt{sw} abbreviate ``store word.'' The source register is the first operand.
Neither source register is erased or advanced by the store.

The operand \texttt{0(\$t0)} combines a constant byte offset with a base register.
The instruction adds zero to the address in \texttt{\$t0} to obtain the first
address it will access, called the \emph{effective address}. An operand of
\texttt{4(\$t0)} refers four bytes farther on. The offset counts bytes; it is not
multiplied by four because the instruction stores a word.

To read the stored value, \texttt{lw \$t2, 0(\$t0)} loads the four-byte word
into \texttt{\$t2}, sign-extending it to the 64-bit register width. Here the first
operand is a destination. The load does not change memory or the base register.
The distinction between load and store is essential: a line that names a register
and an address might change either one, depending on its mnemonic.

Our color has its top word bit clear, so sign extension fills the upper register
bits with zeros. A word such as \texttt{0xFFFFFFFF} would instead load with
\texttt{lw} as \texttt{0xFFFFFFFFFFFFFFFF}, or signed $-1$. The extension rules
from Chapter~\ref{revision:integer-representation} apply when a value comes from
memory as well as when it comes from another register.

\section{Four bytes make one color word}

The MARS display interprets the low 24 bits of a word as red, green and blue
components in the form \texttt{0x00RRGGBB}. Each component occupies eight bits and
ranges from 0 through 255. The high byte is ignored by this RGB display; it is not
an alpha value controlling transparency.

For \texttt{0x00336699}, red is \texttt{0x33}=51, green is
\texttt{0x66}=102 and blue is \texttt{0x99}=153. These positions describe the
numerical word. They do not yet tell us which byte occupies the lowest address.
That depends on \emph{byte order}, also called endianness.

This example uses MARS's little-endian memory. The least significant byte of a
word is stored at the lowest address. The initial color therefore has the
following arrangement. Every byte in the table is written in hexadecimal.

\begin{center}
\begin{tabular}{lll}
\toprule
Address & Byte & Meaning in the color word \\
\midrule
\texttt{0x10010000} & \texttt{99} & Blue \\
\texttt{0x10010001} & \texttt{66} & Green \\
\texttt{0x10010002} & \texttt{33} & Red \\
\texttt{0x10010003} & \texttt{00} & Ignored high byte \\
\bottomrule
\end{tabular}
\end{center}

A word view may show \texttt{00336699} while a byte view at increasing addresses
shows \texttt{99 66 33 00}. Both describe the same memory. Endianness changes the
relationship between byte addresses and significance within a multibyte value.
It does not reverse the bits within each byte or change how we write a register's
hexadecimal value.

\section{Changing exactly one byte}

The instruction \texttt{sb \$t5, 0(\$t0)} stores only the low eight bits of
\texttt{\$t5}. With 255 in that source register, it writes \texttt{FF} at the
lowest-addressed byte of our color. The other three bytes remain unchanged.
The new sequence is \texttt{FF 66 33 00}, representing the word
\texttt{0x003366FF}. Only blue has changed, from 153 to 255.

Store width and register width are separate decisions. Even if the source held
\texttt{0x1234FF}, this byte store would still write only \texttt{FF}. Replacing
\texttt{sb} with \texttt{sw} while the source holds 255 would write
\texttt{FF 00 00 00}. That would erase the red and green components. The correct
first byte alone is insufficient evidence that a word store is appropriate.

We will also load the original blue byte twice. \texttt{lbu} loads one byte and
zero-extends it, giving 153 for \texttt{0x99}. \texttt{lb} sign-extends the same
byte, giving $-103$. The two loads leave memory unchanged. For a color component,
the unsigned interpretation is useful; the signed load is included to connect
this observation with the preceding chapter's two interpretations of one pattern.

\section{The complete program}

The header is a comment consumed by the MARS plugin. It asks Bitmap Display for
a 64-by-64 display with 32-by-32 units, producing two logical columns and two rows.
The first word represents the upper-left cell. The comment is not a processor
instruction. The final eight reserved bytes provide space for the remaining two
cells; the second word is deliberately used as a neighboring-memory check.

\MARSCodeLabel{Store a color, change its blue byte, and reload}
\begin{minted}{mars-idea}
# mars:tool bitmap width=64 height=64 unitWidth=32 unitHeight=32 base=static_data
.data
.align 2
pixel:    .word 0
neighbor: .word 0x13579bdf
          .space 8

.text
.globl main
main:
    la $t0, pixel
    li $t1, 0x00336699
store_color:
    sw $t1, 0($t0)
inspect_initial:
    lw $t2, 0($t0)
    lbu $t3, 0($t0)
    lb $t4, 0($t0)
    li $t5, 0xff
change_blue:
    sb $t5, 0($t0)
inspect_final:
    lw $t6, 0($t0)
    lw $t7, 4($t0)
    li $v0, 10
    syscall
\end{minted}

Save this source as \texttt{memory-pixel.asm}. Run with the MIPS64 Release 6
target and Default memory layout. The program prints nothing to the console.
Its results are memory and register changes. Open the \emph{MARS Tools} window and select
\emph{Bitmap Display} to inspect the connected run's color. Use Debug to examine the
intermediate states rather than expecting two separate colors to remain visible
after a fast Run.

\texttt{.align 2} aligns the following data to a $2^2=4$ byte boundary.
\texttt{.space 8} reserves eight bytes. The example places four-byte objects on
nonoverlapping word boundaries. On this Release 6 target, ordinary word loads
and stores also support unaligned addresses. That support does not make every
offset correct: a word access beginning one byte later still touches four bytes,
possibly including part of a neighboring object.

\section{Follow the memory changes}

At \texttt{store\_color}, check the prepared address and color. Step over the
store, inspect the four bytes and confirm that \texttt{\$t0} is unchanged.

Continue to \texttt{change\_blue}. The first load has copied the original word
into \texttt{\$t2}; the byte loads have placed 153 in \texttt{\$t3} and
$-103$ in \texttt{\$t4}. Step over \texttt{sb} and compare memory with those
registers. Memory now contains the new blue byte, but the earlier register copies
retain their old values. Loading creates a snapshot of a value at that execution
boundary, not a live connection between a register and memory.

Step over the two final loads. \texttt{\$t6} should hold
\texttt{0x00000000003366FF}, and \texttt{\$t7} should still hold
\texttt{0x0000000013579BDF}. Distinguish the register containing the address,
the register containing the earlier word, and the register containing the current
word. Seeing three different values is expected because those registers have
different roles.

\section{Find a wrong address or a wrong width}

In a working copy, change the byte store's offset from zero to one. Predict the
four resulting bytes before running. On this little-endian target, the green
component becomes 255 and the word becomes \texttt{0x0033FF99}. The program still
exits normally. A successful exit does not establish that the right component
changed.

Restore the offset, then replace the byte store with a word store. The result
becomes \texttt{0x000000FF}. Red and green have been cleared even though the
blue byte is correct. Investigate by writing down the effective address and the
set of byte locations touched by the store. These two facts distinguish the
address error from the width error.

The permitted drawing operation owns the four bytes at \texttt{pixel}; it has no
reason to modify \texttt{neighbor}. The adjacent word is a \emph{guard value}:
a recognizable pattern that helps reveal an out-of-bounds write. Its unchanged
value is useful evidence for this example, although one guard cannot prove that
an arbitrary program never writes anywhere else.

\section{Write a related memory operation}

Starting from the complete program, change the red component to \texttt{0xAA}
while preserving green and blue. Predict the new word and every byte before
choosing an offset. Keep the guard unchanged, and use an unsigned byte load to
read back the component you changed. Explain why that byte load and a word load
report different-sized values from the same object.

Then write a version that copies the original four-byte color to a second,
separately declared destination before changing the first pixel. Give the copy its
own label and retain an additional guard. After the update, the original should
contain the new color while the copy retains \texttt{0x00336699}. Identify which
instructions read memory and which write it. A load into a register alone does
not create another memory object.

